\subsection{(\S2) Once More with Feeling}

\begin{enumerate}

    \item \fbox{\it feeling is rational} I, for instance, have trained myself,
        when I lecture on the theory of psychological types, to teach the
        difference I see between feelings and feeling. \ul{Feeling, I tell my
        audiences, is the function that sorts out feelings.} This is my way of
        differentiating feeling from emotion, the first step in that conceptual
        leap Jung made to describe feeling as a rational function, which Jo
        Wheelwright was fond of describing in his training seminars as “a
        landmark in the history of thought.”

    \item \fbox{\it analytic persona} The development of the analytic attitude
        in the microcosm of an analyst's office is not unlike the development
        of the religious attitude that von Franz implies will take a very long
        time to restore in the outer world. What this analytic attitude finally
        might look like I have only had glimpses of, but it is definitely an
        eros tinged by logos that accepts the distances between people and yet
        enables sane-making greeting and meeting and that is filled not only
        with affection for the person I am working with, but also love for the
        work itself. At its best, it unites extraverted and introverted feeling
        into something akin to appreciation—appreciation of the shine that
        comes from someone else's shoes when you get very close to what it is
        like to be in them and have done what you can to make them nicer for
        the person who has to wear them.

\end{enumerate}
